% Options for packages loaded elsewhere
\PassOptionsToPackage{unicode}{hyperref}
\PassOptionsToPackage{hyphens}{url}
\PassOptionsToPackage{dvipsnames,svgnames,x11names}{xcolor}
\documentclass[
  11pt,
]{article}
\usepackage{xcolor}
\usepackage[margin=1in]{geometry}
\usepackage{amsmath,amssymb}
\setcounter{secnumdepth}{-\maxdimen} % remove section numbering
\usepackage{iftex}
\ifPDFTeX
  \usepackage[T1]{fontenc}
  \usepackage[utf8]{inputenc}
  \usepackage{textcomp} % provide euro and other symbols
\else % if luatex or xetex
  \usepackage{unicode-math} % this also loads fontspec
  \defaultfontfeatures{Scale=MatchLowercase}
  \defaultfontfeatures[\rmfamily]{Ligatures=TeX,Scale=1}
\fi
\usepackage{lmodern}
\ifPDFTeX\else
  % xetex/luatex font selection
\fi
% Use upquote if available, for straight quotes in verbatim environments
\IfFileExists{upquote.sty}{\usepackage{upquote}}{}
\IfFileExists{microtype.sty}{% use microtype if available
  \usepackage[]{microtype}
  \UseMicrotypeSet[protrusion]{basicmath} % disable protrusion for tt fonts
}{}
\makeatletter
\@ifundefined{KOMAClassName}{% if non-KOMA class
  \IfFileExists{parskip.sty}{%
    \usepackage{parskip}
  }{% else
    \setlength{\parindent}{0pt}
    \setlength{\parskip}{6pt plus 2pt minus 1pt}}
}{% if KOMA class
  \KOMAoptions{parskip=half}}
\makeatother
\usepackage{longtable,booktabs,array}
\newcounter{none} % for unnumbered tables
\usepackage{calc} % for calculating minipage widths
% Correct order of tables after \paragraph or \subparagraph
\usepackage{etoolbox}
\makeatletter
\patchcmd\longtable{\par}{\if@noskipsec\mbox{}\fi\par}{}{}
\makeatother
% Allow footnotes in longtable head/foot
\IfFileExists{footnotehyper.sty}{\usepackage{footnotehyper}}{\usepackage{footnote}}
\makesavenoteenv{longtable}
\setlength{\emergencystretch}{3em} % prevent overfull lines
\providecommand{\tightlist}{%
  \setlength{\itemsep}{0pt}\setlength{\parskip}{0pt}}
\usepackage{bookmark}
\IfFileExists{xurl.sty}{\usepackage{xurl}}{} % add URL line breaks if available
\urlstyle{same}
\hypersetup{
  pdftitle={Vector-like states of the 27 at two thresholds: a dated prediction αs(MZ) = 0.1178 and the Higgs mass from a vanishing quartic at MPl/√2},
  pdfauthor={Matthew Gelfman (ORCID 0009-0000-6402-2258); Melissa Behrends-Gelfman (ORCID 0009-0001-6281-1947)},
  colorlinks=true,
  linkcolor={Maroon},
  filecolor={Maroon},
  citecolor={Blue},
  urlcolor={Blue},
  pdfcreator={LaTeX via pandoc}}

\usepackage{amssymb}
\usepackage{newunicodechar}
\newunicodechar{±}{\ensuremath{\pm}}
\newunicodechar{²}{\ensuremath{^{2}}}
\newunicodechar{³}{\ensuremath{^{3}}}
\newunicodechar{¹}{\ensuremath{^{1}}}
\newunicodechar{×}{\ensuremath{\times}}
\newunicodechar{Δ}{\ensuremath{\Delta}}
\newunicodechar{α}{\ensuremath{\alpha}}
\newunicodechar{θ}{\ensuremath{\theta}}
\newunicodechar{λ}{\ensuremath{\lambda}}
\newunicodechar{μ}{\ensuremath{\mu}}
\newunicodechar{π}{\ensuremath{\pi}}
\newunicodechar{σ}{\ensuremath{\sigma}}
\newunicodechar{⁰}{\ensuremath{^{0}}}
\newunicodechar{⁴}{\ensuremath{^{4}}}
\newunicodechar{⁶}{\ensuremath{^{6}}}
\newunicodechar{⁸}{\ensuremath{^{8}}}
\newunicodechar{⁺}{\ensuremath{^{+}}}
\newunicodechar{⁻}{\ensuremath{^{-}}}
\newunicodechar{₁}{\ensuremath{_{1}}}
\newunicodechar{₂}{\ensuremath{_{2}}}
\newunicodechar{₃}{\ensuremath{_{3}}}
\newunicodechar{₆}{\ensuremath{_{6}}}
\newunicodechar{→}{\ensuremath{\to}}
\newunicodechar{−}{\ensuremath{-}}
\newunicodechar{√}{\ensuremath{\surd}}
\newunicodechar{≈}{\ensuremath{\approx}}
\title{Vector-like states of the 27 at two thresholds: a dated
prediction \ensuremath{\alpha}\textsubscript{s}(M\textsubscript{Z}) = 0.1178 and the Higgs
mass from a vanishing quartic at M\textsubscript{Pl}/\ensuremath{\surd}2}
\author{Matthew Gelfman (ORCID 0009-0000-6402-2258) \and Melissa
Behrends-Gelfman (ORCID 0009-0001-6281-1947)}
\date{25 November 2026}

\begin{document}
\maketitle

\subsection{Abstract}

Each 27 of E\ensuremath{_{6}} contains, besides one Standard Model generation and a
neutral singlet, a vector-like colour triplet D(3,1,\ensuremath{-}1/3) and a
vector-like weak doublet L(1,2,\ensuremath{-}1/2) (Gürsey, Ramond and Sikivie 1976).
We place three generations of each at two thresholds. The D states sit
at M\textsubscript{D} = 4.6 \ensuremath{\times} 10\ensuremath{^{1}}\ensuremath{^{3}} GeV, and the doublets at
M\textsubscript{L} = M\textsubscript{D}\ensuremath{^{2}}/E\textsubscript{g} = 2.4 \ensuremath{\times} 10\ensuremath{^{8}}
GeV, where E\textsubscript{g} = M\textsubscript{Pl}/\ensuremath{\surd}2. With the
measured electroweak inputs, the three gauge couplings then cross at
M\textsubscript{D} to 0.02 in 1/\ensuremath{\alpha} at two loops. Closing the crossing at
three loops requires \ensuremath{\alpha}\textsubscript{s}(M\textsubscript{Z}) = 0.1178.
That is 0.2\ensuremath{\sigma} from the world average (0.1180 \ensuremath{\pm} 0.0009) and 1.7\ensuremath{\sigma} below the
most precise lattice determination (0.11873 \ensuremath{\pm} 0.00056; Dalla Brida et
al.~2025), so lattice results at about 0.3\% will decide it.

The same states change the running of the Higgs quartic. With the
boundary condition \ensuremath{\lambda}(E\textsubscript{g}) = 0 of Shaposhnikov and
Wetterich (2010), the Higgs mass is 127.2 \ensuremath{\pm} 1.5 GeV, against 125.13 \ensuremath{\pm}
0.11 GeV (+1.4\ensuremath{\sigma}). The Standard Model alone gives 128.8 \ensuremath{\pm} 1.5 GeV
(+2.6\ensuremath{\sigma}). Read the other way, the rule needs a top pole mass of 171.6
GeV, 3.7\ensuremath{\sigma} below today's 172.60 \ensuremath{\pm} 0.27 GeV.

Three choices were made with the data in view: E\textsubscript{g} (the
best of five equally natural readings), the M\textsubscript{L} rule, and
equality of the couplings at M\textsubscript{D}. We state each one,
together with what would falsify the construction.

\subsection{1. Field content and thresholds}

Each generation adds one Dirac colour triplet D with hypercharge \ensuremath{-}1/3
and one Dirac weak doublet L with hypercharge \ensuremath{-}1/2. No coloured scalars
are added, so baryon number is conserved at the renormalisable level,
consistent with the proton-lifetime limit above 2.4 \ensuremath{\times} 10\ensuremath{^{3}}\ensuremath{^{4}} years
(Takenaka et al.~2020). The D states mix only with d\textsubscript{R}.
Their mixing with left-handed quarks is suppressed by
(v/M\textsubscript{D})\ensuremath{^{2}} \ensuremath{\approx} 10\ensuremath{^{-}}\ensuremath{^{2}}\ensuremath{^{3}}, so first-row CKM unitarity holds to
that level.

In GUT normalisation, each Dirac D contributes \ensuremath{\Delta}b = (4/15, 0, 2/3) at
one loop, and each Dirac L contributes \ensuremath{\Delta}b = (2/5, 2/3, 0). The full
two-loop matrices are used (Machacek and Vaughn 1984; Luo, Wang and Xiao
2003).

\subsection{\texorpdfstring{2. The crossing and
\ensuremath{\alpha}\textsubscript{s}}{2. The crossing and alphas}}

With sin\ensuremath{^{2}}\ensuremath{\theta}\textsubscript{W} = 0.23122, 1/\ensuremath{\alpha}(M\textsubscript{Z}) = 127.951
and \ensuremath{\alpha}\textsubscript{s}(M\textsubscript{Z}) = 0.1178, the inverse
couplings (GUT-normalised) run as follows:

{\def\LTcaptype{none} % do not increment counter
\begin{longtable}[]{@{}llll@{}}
\toprule\noalign{}
\ensuremath{\mu} (GeV) & 1/\ensuremath{\alpha}\ensuremath{_{1}} & 1/\ensuremath{\alpha}\ensuremath{_{2}} & 1/\ensuremath{\alpha}\ensuremath{_{3}} \\
\midrule\noalign{}
\endhead
\bottomrule\noalign{}
\endlastfoot
2.4 \ensuremath{\times} 10\ensuremath{^{8}} & 49.31 & 36.92 & 25.33 \\
4.6 \ensuremath{\times} 10\ensuremath{^{1}}\ensuremath{^{3}} & 38.97 & 38.99 & 38.99 \\
8.63 \ensuremath{\times} 10\ensuremath{^{1}}\ensuremath{^{8}} & 27.10 & 41.08 & 48.59 \\
\end{longtable}
}

The couplings cross at M\textsubscript{D} and separate again above it.
No placement of the two thresholds makes them equal at
E\textsubscript{g}: the smallest spread found is 16.8 in 1/\ensuremath{\alpha}. The
construction therefore takes equality at M\textsubscript{D}, where each
27 becomes complete, as a boundary condition, because there is no
unified group to enforce it.

The prediction moves with the choice of E\textsubscript{g}. With the
unreduced Planck energy it is \ensuremath{\alpha}\textsubscript{s} = 0.1196, which is not
the value claimed here. Because E\textsubscript{g} was chosen as the
best of five readings, the prediction is tested only by measurements
made after it was fixed (17 November 2026).

\subsection{3. The Higgs mass}

The two-loop running starts from the NNLO matching of Buttazzo et
al.~(2013), and the Standard Model's three-loop shift is added as an
estimate. The uncertainty includes the top pole mass (\ensuremath{\pm}0.27 GeV plus 0.5
GeV for the pole/Monte-Carlo ambiguity), \ensuremath{\alpha}\textsubscript{s} and missing
higher orders. The condition \ensuremath{\lambda}(E\textsubscript{g}) = 0 then gives:

\begin{itemize}
\tightlist
\item
  Standard Model only: m\textsubscript{H} = 128.8 \ensuremath{\pm} 1.5 GeV.
\item
  With the D and L states: m\textsubscript{H} = 127.2 \ensuremath{\pm} 1.5 GeV. Without
  the 0.5 GeV allowance the tension is 2.1\ensuremath{\sigma}.
\end{itemize}

The pass line, fixed before running, was 2\ensuremath{\sigma}. The result passes only with
the allowance.

\subsection{4. What the construction does not do}

The rule M\textsubscript{L} = M\textsubscript{D}\ensuremath{^{2}}/E\textsubscript{g} was
found after the needed value (2.2 \ensuremath{\times} 10\ensuremath{^{8}} GeV) was known. A companion note
shows that no flavour-blind discrete charge, of any order, can enforce
it without a gauge anomaly. The rule is therefore an assumption: the
doublets couple to the field that gives D its mass only through
gravitational-strength terms.

If the universe reheated above M\textsubscript{L}, the doublets must
decay through a small mixing with the lepton doublets, of order 10\ensuremath{^{-}}\ensuremath{^{1}}\ensuremath{^{6}} or
larger depending on the reheating temperature. That adds one coupling.

\subsection{5. Falsifiers}

The construction fails if any of the following is observed:

\begin{itemize}
\tightlist
\item
  \ensuremath{\alpha}\textsubscript{s}(M\textsubscript{Z}) more than 3\ensuremath{\sigma} from 0.1178 at
  about 0.3\% precision;
\item
  a top mass that puts the Higgs condition more than 3\ensuremath{\sigma} off;
\item
  any proton decay (today's limit: Takenaka et al.~2020);
\item
  a confirmed shortfall in first-row CKM unitarity (today 0.9984 \ensuremath{\pm}
  0.0007, a 2.3\ensuremath{\sigma} shortfall);
\item
  coloured scalars or extra vector-like matter below M\textsubscript{L}.
\end{itemize}

\subsection{Data and code}

\texttt{analysis/leak\_establish\_three\_loop\_2026-11-16.py},
\texttt{analysis/crossing\_vs\_grain\_2026-11-18.py},
\texttt{analysis/higgs\_at\_the\_grain\_2026-11-17.py} and
\texttt{analysis/higgs\_three\_loop\_check\_2026-11-18.py} in the TRPG
code archive, doi:10.5281/zenodo.24329239. This result is one component
of the framework in doi:10.5281/zenodo.24330763.

\subsection{Acknowledgements}

Due to medical necessity, AI language-model assistants were used, under
the authors' direction, to draft and edit text from the authors' ideas,
to carry out and check calculations, and to search the literature; every
result was checked by rerunning the published code. The authors reviewed
all content and take full responsibility for it.

\subsection{References}

\begin{itemize}
\tightlist
\item
  Buttazzo, D., Degrassi, G., Giardino, P. P., Giudice, G. F., Sala, F.,
  Salvio, A. and Strumia, A. (2013). Investigating the near-criticality
  of the Higgs boson. \emph{JHEP} 12, 089.
\item
  Dalla Brida, M., Höllwieser, R., Knechtli, F., Korzec, T., Ramos, A.,
  Sint, S. and Sommer, R. (ALPHA) (2025). The strength of the
  interaction between quarks and gluons. arXiv:2501.06633.
\item
  Gürsey, F., Ramond, P. and Sikivie, P. (1976). A universal gauge
  theory model based on E6. \emph{Phys. Lett. B} 60, 177.
\item
  Luo, M., Wang, H. and Xiao, Y. (2003). Two-loop renormalization group
  equations in general gauge field theories. \emph{Phys. Rev.~D} 67,
  065019.
\item
  Machacek, M. E. and Vaughn, M. T. (1984). Two-loop renormalization
  group equations in a general quantum field theory. II. Yukawa
  couplings. \emph{Nucl. Phys. B} 236, 221.
\item
  Shaposhnikov, M. and Wetterich, C. (2010). Asymptotic safety of
  gravity and the Higgs boson mass. \emph{Phys. Lett. B} 683, 196.
\item
  Takenaka, A. et al.~(Super-Kamiokande Collaboration) (2020). Search
  for proton decay via p \ensuremath{\to} e\ensuremath{^{+}}\ensuremath{\pi}\ensuremath{^{0}} and p \ensuremath{\to} \ensuremath{\mu}\ensuremath{^{+}}\ensuremath{\pi}\ensuremath{^{0}} with an enlarged fiducial
  volume in Super-Kamiokande I--IV. \emph{Phys. Rev.~D} 102, 112011.
\end{itemize}

\end{document}
